\documentclass[../cours_spe_trigo.tex]{subfiles}
\begin{document}
\setcounter{chapter}{18}
\chapter{Synthèse finale : savoir reconstruire plutôt que retenir}
\begin{tcolorbox}[enhanced,colback=bleuPale,colframe=bleu,boxrule=.8pt,arc=1mm]
\begin{align*}
\pi&=C/D\quad\text{(après preuve de constance du rapport)},\\
C&=2\pi R,\qquad \mathcal A_{\mathrm{disque}}=\pi R^2,\\
\theta_{\mathrm{rad}}&=\ell/R,\qquad 180^\circ=\pi\,\mathrm{rad},\\
\ell&=R\theta,\qquad \mathcal A_{\mathrm{secteur}}=\tfrac12 R^2\theta\quad (0\le\theta\le2\pi),\\
M(\theta)&=(\cos\theta,\sin\theta),\\
\cos^2\theta+\sin^2\theta&=1,\qquad \tan\theta=\sin\theta/\cos\theta\ \text{si }\cos\theta\ne0,\\
\cos(a+b)&=\cos a\cos b-\sin a\sin b,\\
\sin(a+b)&=\sin a\cos b+\cos a\sin b,\\
\cos2x&=\cos^2x-\sin^2x,\qquad \sin2x=2\sin x\cos x,\\
(\sin x)'&=\cos x,\qquad(\cos x)'=-\sin x,\\
\lim_{x\to0}\sin x/x&=1\quad\text{(en radians)},\\
\vec u\cdot\vec v&=\|\vec u\|\|\vec v\|\cos\gamma.
\end{align*}
\end{tcolorbox}
\textbf{Le fil directeur.} Les polygones expliquent la longueur du cercle ; les homothéties expliquent $\pi$ ; les radians relient angle et longueur ; les triangles expliquent les rapports ; les coordonnées donnent le cercle unité ; les rotations fournissent les additions ; les additions permettent les équations et les dérivées ; l'analyse mène aux oscillations et aux nombres complexes.

\textbf{Question-test.} Face à une identité nouvelle, demander : \emph{vient-elle d'une définition, de Pythagore, d'une symétrie, d'une rotation, d'une limite ou d'une règle algébrique ?} Si l'on sait répondre à cette question, on possède les fondements du sujet.


\end{document}
